\section{What each component contributes}
\label{sec:ablation}

Two reviews observe that Table~3 is cumulative, that it does not say which components each row contains, and that its last row adds a warp, more matching, verification features and a second fusion at once. \Cref{sec:inventory} gives the inventory of every row and corrects two of them. \Cref{sec:controls} isolates the warp from the matching budget and from the learned fusion with dedicated controls, and \cref{sec:iteration} tests whether the refinement can be iterated without moving.

\subsection{Component inventory of Table~3}
\label{sec:inventory}

\begin{table*}[t]
\centering
\caption{\textbf{Component inventory} on the 6{,}364 test pairs. All columns but the last are in the camera frame and involve no absolute heading. Pairs whose pose cannot be solved predict zero motion (``fail''). Angular errors are mean / median. The last column is the endpoint success in ENU with the learned absolute heading.}
\label{tab:inventory}
\footnotesize
\setlength{\tabcolsep}{4pt}
\resizebox{\textwidth}{!}{%
\begin{tabular}{@{}lcccccccccc@{}}
\toprule
& fail & $d$ MAE & \SRd@1\,m & \multicolumn{2}{c}{EPE (m)} & \SRe@1\,m & $\dth$ error ($^\circ$) & $\dpsi$ error ($^\circ$) & Dir. med. & ENU \\
Configuration & (\%) & (m) & (\%) & mean & med. & (\%) & mean / med. & mean / med. & ($^\circ$) & \SRe@1\,m \\
\midrule
\multicolumn{11}{@{}l}{\emph{Learned priors only (no correspondences)}}\\
P\quad MLP on pooled MASt3R statistics & 0.0 & 3.15 & 21.5 & 6.73 & 6.08 & 0.6 & 1.82 / 1.140 & 1.86 / 1.350 & 14.98 & 0.1 \\
\midrule\multicolumn{11}{@{}l}{\emph{Two-view geometry only: no learned prior, no fusion; scale from the flight-height prior}}\\
G0\quad Frozen SuperGlue, one pass, E & 53.8 & 14.71 & 22.8 & 16.12 & 13.69 & 17.5 & 68.06 / 73.467 & 12.56 / 6.455 & 90.00 & -- \\
G1\quad Rotation-aware matching, E & 0.7 & 2.97 & 58.3 & 4.78 & 1.11 & 46.1 & 1.43 / 0.032 & 1.85 / 0.180 & 1.65 & -- \\
G3\quad \quad + zoom re-matching, E & 1.1 & 2.22 & 64.6 & 3.48 & 0.91 & 54.1 & 0.99 / 0.026 & 1.33 / 0.136 & 1.23 & -- \\
G5\quad \quad + imagination re-matching, E & 0.5 & 1.19 & 80.3 & 1.70 & 0.55 & 77.2 & 0.59 / 0.007 & 0.59 / 0.036 & 0.36 & -- \\
\midrule\multicolumn{11}{@{}l}{\emph{Geometry with learned priors in the pose refinement (\cref{eq:map}); scale from the flight-height prior; no fusion}}\\
G2\quad Rotation-aware matching, refined & 0.7 & 2.11 & 60.2 & 3.55 & 1.07 & 47.2 & 1.26 / 0.032 & 1.24 / 0.170 & 1.57 & -- \\
G4\quad \quad + zoom re-matching, refined & 1.0 & 1.59 & 65.8 & 2.57 & 0.92 & 53.8 & 0.92 / 0.027 & 0.96 / 0.135 & 1.23 & -- \\
G6\quad \quad + imagination re-matching, refined & 0.5 & 0.98 & 80.6 & 1.37 & 0.56 & 77.4 & 0.57 / 0.007 & 0.47 / 0.037 & 0.38 & -- \\
\midrule\multicolumn{11}{@{}l}{\emph{Same geometry, scale from the learned distance regression instead of the flight-height prior}}\\
G2$'$\quad Rotation-aware matching, refined & 0.7 & 2.79 & 24.9 & 4.87 & 2.93 & 16.0 & 1.26 / 0.032 & 1.24 / 0.170 & 1.57 & -- \\
G6$'$\quad Imagination re-matching, refined & 0.5 & 2.76 & 25.0 & 3.32 & 2.28 & 22.9 & 0.57 / 0.007 & 0.47 / 0.037 & 0.38 & -- \\
\midrule\multicolumn{11}{@{}l}{\emph{Planar pose model: homography decomposition instead of the essential matrix; scale from the flight height; no learning}}\\
H1\quad Rotation-aware matching, H & 14.8 & 21.12 & 16.7 & 26.02 & 7.54 & 2.7 & 14.18 / 0.287 & 8.17 / 1.790 & 18.48 & -- \\
H3\quad \quad + zoom re-matching, H & 1.4 & 7.03 & 24.8 & 11.36 & 4.82 & 5.0 & 1.63 / 0.161 & 3.24 / 1.030 & 11.27 & -- \\
H5\quad \quad + imagination re-matching, H & 0.8 & 5.00 & 27.7 & 8.64 & 4.17 & 6.3 & 1.11 / 0.132 & 2.50 / 0.820 & 9.31 & -- \\
\midrule\multicolumn{11}{@{}l}{\emph{Reliability-weighted fusion of all candidates and scale cues (the rows of Table~3)}}\\
F1\quad Rotation-aware pass only & 0.0 & 1.14 & 65.1 & 2.43 & 1.05 & 48.0 & 0.19 / 0.030 & 0.64 / 0.162 & 1.63 & 35.2 \\
F2\quad + zoom re-matching (Stage I) & 0.0 & 0.92 & 73.1 & 1.71 & 0.82 & 59.9 & 0.12 / 0.021 & 0.44 / 0.105 & 1.04 & 44.3 \\
F3\quad + geometric imagination (Stage II) & 0.0 & 0.68 & 83.3 & 1.01 & 0.53 & 78.7 & 0.07 / 0.006 & 0.20 / 0.028 & 0.35 & 61.9 \\
\bottomrule
\end{tabular}}
\end{table*}


\Cref{tab:inventory} recomputes the chain on all \nTest{} test pairs, in the camera frame, so that no row depends on the absolute heading. It separates three kinds of configuration that Table~3 mixes: the learned priors alone (P), two-view geometry without any fusion (G0--G6), and the fused estimates (F1--F3). In the G rows the pose comes from one matching pass and the length of the translation from the flight-height cue of \cref{sec:cands} with its single constant; nothing else is learned in G0, G1, G3 and G5. The H rows replace the essential matrix by the decomposition of a plane-induced homography, the classical estimator for planar scenes.

\noindent\textbf{Rows of Table~3.}
\begin{itemize}[leftmargin=*,itemsep=1pt,topsep=2pt]
\item \emph{``Frozen SuperGlue''} (5.14\,m, 11.7\%). This row is mislabelled. It is the regression from pooled frozen-MASt3R statistics with the original heads of RASR~\cite{liang2026rasr}, evaluated in ENU; it uses no correspondences. With the camera-frame target the same features give row~P. SuperGlue used literally, one pass without rotation handling, is row~G0: it finds no pose for \gZeroFail\% of the pairs, whose headings differ arbitrarily, and reaches \gZeroSRd\% \SRd.
\item \emph{``+ rotation-aware matching''} (4.80\,m, 12.3\%). This is the fused estimate F1, but scored in ENU \emph{without} an absolute heading: $\Delta z$, $\dth$ and $\dpsi$ come from the fusion, while $(\Delta x,\Delta y)$ cannot be formed from $(\ell_x,\ell_y)$ and are taken from the world-frame regression of the prior heads, which is no better than predicting zero. The 4.80\,m therefore measure the missing horizontal component, not the quality of the geometry: in the camera frame the same estimate has $d{=}$\fOneD\,m and \SRd{}${=}$\fOneSRd\%.
\item \emph{``+ absolute heading''} (printed: 2.25\,m, 35.1\%). With a heading, $(\ell_x,\ell_y)$ is rotated into ENU. This cannot change the range error, as one reviewer notes, and it does not: the row has $d{=}$\fOneD\,m and \SRd{}${=}$\fOneSRd\% with or without the heading. The two values printed in Table~3 are wrong. The printed success coincides with the \emph{endpoint} success of this configuration in ENU (\fOneEnuSRe\%), not with the range-based success of the other rows, and we could not reproduce the printed 2.25\,m. The sentence of Sec.~5.2 that attributes a reduction of the distance error from 4.80\,m to 2.25\,m to the heading should be disregarded.
\item \emph{``+ zoom re-matching (Stage~I)''} and \emph{``+ geometric imagination (Stage~II)''} are F2 and F3 and are correct as printed.
\end{itemize}
With the corrected third row, the zoom pass adds 8.0 points of \SRd{} (\fOneSRd\% to \SiCamSRd\%) and not 38.0, and the imagination-refined stage adds \gainSRd{} points; the refinement is the larger of the two steps, also under the endpoint definition (\fOneSRe\%, \SiCamSRe\%, \SiiCamSRe\%). All fused rows F1--F3 contain the same priors, scale cues and fusion; they differ only in the matching passes that feed candidates to it.

\noindent\textbf{What the inventory shows.}
(i) \emph{Geometry carries the estimate.} The priors alone reach \priSRd\% \SRd{}. A calibrated two-view pipeline built from the rotation-aware matcher, with no learned component other than the flight-height constant, reaches \gOneSRd\% (G1), and \gFourSRd\% with the zoom pass and the prior-regularised pose (G4). These are the appropriate references for DreamNav, well above the 11.3\% of the best policy in Table~2; \cref{sec:geometric} extends them to other matchers.
(ii) \emph{Re-matching after the plane-induced warp is the largest single step, and it does not need the second fusion.} Replacing the zoom matches by the imagined match set lifts the geometry-only estimate from \gFourSRd\% to \gSixSRd\% in \SRd{} and from \gFourSRe\% to \gSixSRe\% in \SRe{} (G4 to G6), with no fusion involved.
(iii) \emph{The fusion mainly removes failures.} Relative to the best geometry-only row of each stage, the fusion adds 4.9, 7.3 and 2.7 points of \SRd{} (F1--F3 against G2, G4, G6). Its effect is larger on the means, which failures dominate: from G6 to F3 the mean heading error falls from \gSixTheta$^\circ$ to 0.07$^\circ$, the mean tilt error from \gSixPsi$^\circ$ to 0.20$^\circ$, and the mean range error from \gSixD\,m to 0.68\,m.
(iv) \emph{Metric scale comes from the flight-height prior.} When the length of the translation is taken from the learned distance regression instead, \SRd{} stays near 25\% whatever the geometry (G2$'$, G6$'$). The calibrated depth cue of MASt3R and the direct regression are useful to the fusion as fallbacks, but the cue that makes a metre-level range possible on PairUAV is the regularity of the flight height, a property of the dataset that DreamNav and the geometric baselines of \cref{sec:geometric} share. We come back to this dependence in \cref{sec:planarity,sec:real,sec:corrections}.
(v) \emph{The plane is not a sufficient pose model.} Decomposing a homography fitted to the same matches, with the calibration and the flight height, in place of the essential matrix (rows H1--H5; inlier threshold selected on training pairs) reaches only \hmOneSRd\%, \hmThreeSRd\% and \hmFiveSRd\% \SRd, with median endpoint errors of \hmOneEPEmed, \hmThreeEPEmed{} and \hmFiveEPEmed\,m: the matched points are not coplanar (\cref{sec:planarity}). This is why DreamNav uses the plane only to align the two images and estimates the pose from epipolar geometry.

\subsection{Is it the warp, the extra matching, or the fusion?}
\label{sec:controls}

\begin{table*}[t]
\centering
\caption{\textbf{Controls for the refinement stage} on 3{,}000 test pairs; every fusion is fitted on the same 3{,}000 training pairs. Camera frame. The last two columns are paired differences to the coarse stage in percentage points with 95\% bootstrap intervals over buildings.}
\label{tab:controls}
\footnotesize
\setlength{\tabcolsep}{4pt}
\resizebox{\textwidth}{!}{%
\begin{tabular}{@{}lccccccccc@{}}
\toprule
& $d$ MAE & \SRd@1\,m & \multicolumn{2}{c}{EPE (m)} & \SRe@1\,m & $\dth$ & $\dpsi$ & $\Delta$\SRd & $\Delta$\SRe \\
Variant & (m) & (\%) & mean & med. & (\%) & MAE ($^\circ$) & MAE ($^\circ$) & \vs coarse & \vs coarse \\
\midrule
Coarse stage (input of all rows below) & 0.91 & 73.5 & 1.69 & 0.82 & 60.5 & 0.132 & 0.405 & -- & -- \\
\midrule\multicolumn{10}{@{}l}{\emph{(a) Learned capacity without new correspondences}}\\
Second fusion stage on the coarse candidates & 0.93 & 72.4 & 1.71 & 0.85 & 58.0 & 0.135 & 0.412 & $-1.1$ \ci{-2.1}{-0.2} & $-2.4$ \ci{-3.4}{-1.4} \\
\midrule\multicolumn{10}{@{}l}{\emph{(b) New correspondences, full fusion; what the second matching pass sees}}\\
Rotation-only alignment, $\le 3{\times}3$ tiles (compute-matched) & 0.72 & 82.1 & 1.11 & 0.57 & 76.0 & 0.086 & 0.234 & $+8.6$ \ci{+7.2}{+10.1} & $+15.6$ \ci{+13.8}{+17.5} \\
Plane-induced warp, 1 tile (no extra budget) & 0.79 & 77.5 & 1.37 & 0.72 & 65.8 & 0.105 & 0.359 & $+4.0$ \ci{+2.8}{+5.3} & $+5.3$ \ci{+3.9}{+6.8} \\
Plane-induced warp at $512{\times}288$, 1 tile & 0.93 & 72.3 & 1.72 & 0.87 & 57.1 & 0.132 & 0.415 & $-1.2$ \ci{-2.2}{-0.4} & $-3.4$ \ci{-4.4}{-2.4} \\
Plane-induced warp, $\le 3{\times}3$ tiles = Stage II & 0.70 & 82.7 & 1.03 & 0.54 & 78.3 & 0.073 & 0.207 & $+9.2$ \ci{+7.7}{+10.7} & $+17.8$ \ci{+16.1}{+19.5} \\
\midrule\multicolumn{10}{@{}l}{\emph{(c) Plane-induced warp with $\le 3{\times}3$ tiles; how the re-estimated pose is used}}\\
Re-estimated pose replaces the coarse pose (no fusion) & 0.93 & 73.8 & 1.59 & 0.65 & 69.1 & 0.171 & 0.484 & $+0.3$ \ci{-0.2}{+0.8} & $+8.6$ \ci{+7.4}{+9.8} \\
Uniform candidate weights, no residual & 1.09 & 68.2 & 1.53 & 0.77 & 61.8 & 0.191 & 0.495 & $-5.4$ \ci{-7.5}{-3.5} & $+1.3$ \ci{-0.8}{+3.5} \\
Predicted-error weights, no residual & 0.72 & 82.0 & 1.03 & 0.54 & 77.9 & 0.074 & 0.213 & $+8.5$ \ci{+7.1}{+10.0} & $+17.4$ \ci{+15.8}{+19.2} \\
Full fusion without verification scores & 0.69 & 83.0 & 1.03 & 0.54 & 77.9 & 0.074 & 0.212 & $+9.5$ \ci{+8.1}{+10.9} & $+17.4$ \ci{+15.7}{+19.2} \\
Full fusion, scale kept from the coarse stage & 0.89 & 74.9 & 1.25 & 0.62 & 70.1 & 0.081 & 0.235 & $+1.3$ \ci{+0.6}{+2.2} & $+9.7$ \ci{+8.4}{+10.9} \\
Full fusion = Stage II & 0.70 & 82.7 & 1.03 & 0.54 & 78.3 & 0.073 & 0.207 & $+9.2$ \ci{+7.7}{+10.7} & $+17.8$ \ci{+16.1}{+19.5} \\
\bottomrule
\end{tabular}}
\end{table*}

\begin{table}[t]
\centering
\caption{\textbf{What each second matching pass delivers before any fusion}: tiles matched (at most $3{\times}3$), median number of inliers, share of pairs with fewer than 100 inliers, median displacement between the aligned and the real view, and accuracy of the pose refined with \cref{eq:map}. Unsolved pairs count as $90^\circ$ / $20^\circ$.}
\label{tab:raw_geometry}
\footnotesize
\setlength{\tabcolsep}{2.5pt}
\resizebox{\columnwidth}{!}{%
\begin{tabular}{@{}lcccccccc@{}}
\toprule
& & Inl. & $<$100 & Res. & Dir. ($^\circ$) & Dir. & $\dth$ & $\dpsi$ \\
Second matching pass & Tiles & med. & (\%) & (px) & med. / mean & $<2^\circ$ & mean & mean \\
\midrule
\multicolumn{9}{@{}l}{\emph{All 3{,}000 test pairs of \cref{tab:controls}}}\\
Zoom pass of the coarse stage & 1 & 468 & 4.4 & -- & 1.22 / 8.4 & 63.7 & 1.18 & 0.86 \\
Rotation-only alignment, tiles & 4.9 & 2353 & 1.0 & 182.1 & 0.44 / 5.2 & 85.2 & 0.52 & 0.51 \\
Plane-induced warp, 1 tile & 1.0 & 605 & 2.5 & 7.4 & 0.97 / 6.4 & 70.6 & 0.75 & 0.69 \\
Plane-induced warp at $512{\times}288$ & 1.0 & 119 & 39.5 & 10.2 & 8.58 / 19.9 & 12.5 & 1.86 & 2.15 \\
Plane-induced warp, tiles (ours) & 5.0 & 3221 & 0.7 & 7.3 & 0.39 / 4.1 & 87.7 & 0.50 & 0.55 \\
\midrule\multicolumn{9}{@{}l}{\emph{Subset of 1{,}000 pairs}}\\
Zoom pass of the coarse stage & 1 & 488 & 4.3 & -- & 1.16 / 7.4 & 66.7 & 0.92 & 0.88 \\
Plane-induced warp, tiles (ours) & 5.1 & 3408 & 0.6 & 7.2 & 0.36 / 3.5 & 89.1 & 0.32 & 0.37 \\
Plane-induced warp from the true pose, tiles & 5.1 & 3404 & 0.0 & 6.3 & 0.38 / 3.4 & 89.6 & 0.04 & 0.25 \\
\midrule\multicolumn{9}{@{}l}{\emph{Subset of 1{,}000 pairs}}\\
No pose-guided alignment, tiles & 9.0 & 12 & 71.5 & 665.8 & 89.23 / 65.0 & 16.7 & 37.59 & 14.51 \\
Rotation-only alignment, tiles & 5.0 & 2447 & 1.0 & 182.1 & 0.43 / 4.8 & 86.4 & 0.52 & 0.46 \\
\bottomrule
\end{tabular}}
\end{table}

\begin{table}[t]
\centering
\caption{\textbf{Where the plane-induced part of the warp helps.} Endpoint success (\%) of the coarse stage, of the compute-matched refinement with rotation-only alignment and of the refinement with the plane-induced warp, and their paired difference in points with a 95\% bootstrap interval. Same pairs and protocol as \cref{tab:controls}.}
\label{tab:plane_vs_rot}
\footnotesize
\setlength{\tabcolsep}{3.5pt}
\begin{tabular}{@{}lrcccc@{}}
\toprule
& & & Rot. & Plane & Plane $-$ rot. \\
Bin & Pairs & Coarse & only & warp & (points) \\
\midrule
\multicolumn{6}{@{}l}{\emph{Larger tilt of the pair}}\\
0-15$^\circ$ & 396 & 87.9 & 91.9 & 91.9 & $+0.0$ \ci{-1.8}{+2.0} \\
15-25$^\circ$ & 662 & 76.6 & 85.8 & 84.7 & $-1.1$ \ci{-3.1}{+0.9} \\
25-35$^\circ$ & 894 & 60.5 & 76.7 & 79.1 & $+2.3$ \ci{+0.2}{+4.4} \\
35-45$^\circ$ & 1048 & 39.9 & 63.3 & 68.3 & $+5.1$ \ci{+2.6}{+7.6} \\
\midrule\multicolumn{6}{@{}l}{\emph{Altitude change $|\Delta z|$}}\\
0-5\,m & 478 & 57.1 & 77.4 & 77.6 & $+0.2$ \ci{-2.7}{+3.5} \\
5-15\,m & 949 & 53.7 & 73.6 & 75.2 & $+1.7$ \ci{-0.8}{+4.1} \\
15-25\,m & 1027 & 64.5 & 76.0 & 79.1 & $+3.0$ \ci{+1.2}{+4.8} \\
$\geq$25\,m & 546 & 67.6 & 79.1 & 82.6 & $+3.5$ \ci{+1.0}{+5.9} \\
\midrule\multicolumn{6}{@{}l}{\emph{Relief}}\\
$<$5\% & 241 & 24.9 & 47.3 & 49.8 & $+2.5$ \ci{-3.9}{+8.7} \\
5-10\% & 1803 & 61.7 & 78.6 & 81.0 & $+2.3$ \ci{+0.6}{+3.9} \\
10-20\% & 812 & 69.0 & 80.5 & 82.5 & $+2.0$ \ci{+0.4}{+3.7} \\
$\geq$20\% & 119 & 67.2 & 78.2 & 81.5 & $+3.4$ \ci{-0.8}{+8.7} \\
\bottomrule
\end{tabular}
\end{table}


The controls of \cref{tab:controls} follow the suggestions of the reviewers. They are run on \ctN{} test pairs, and every fusion is fitted with the unchanged protocol on the same \ctTrainRows{} training pairs; with 6{,}000 or 12{,}000 training pairs the complete refinement is within \ctTrainSizeGap{} points of the value reported here on the same test pairs. \Cref{tab:raw_geometry} reports, for the same pairs, what each second matching pass delivers before any learning.

\noindent\textbf{(a) A second fusion stage without new correspondences brings nothing.} Fitting the refinement fusion on the candidates and features of the coarse stage alone changes \SRd{} by \ctFusionOnlydSRd{} points and \SRe{} by \ctFusionOnlydSRe{}. The gain of the refinement stage is not additional learned capacity.

\noindent\textbf{(b) What the second matching pass must see.} All rows of block (b) use the same fusion; they differ in how the two images are aligned before they are matched and in the matching budget.

\emph{Compute-matched pass without the plane-induced warp.} The control requested by the reviewers runs the same tile-wise matching at the original resolution, with the same number of tiles and keypoints, but aligns the two images by the rotation of the coarse pose only: no translation, no plane, no height. It reaches \ctRotSRd\% \SRd{} and \ctRotSRe\% \SRe, \ie \ctRotdSRd{} and \ctRotdSRe{} points over the coarse stage, against \ctFulldSRd{} and \ctFulldSRe{} for the plane-induced warp. The paired difference between the two is \ctRotVsFullSRd{} points of \SRd{} \ci{\ctRotVsFullSRdLo}{\ctRotVsFullSRdHi}, which is not significant, and \ctRotVsFullSRe{} points of \SRe{} \ci{\ctRotVsFullSReLo}{\ctRotVsFullSReHi}. The rotation-aligned pass therefore recovers \ctRotShareD\% and \ctRotShareE\% of the gain of the refinement stage. The plane-induced part of the warp does what Sec.~4.3 says it does for the alignment itself: the median displacement between the two aligned views falls from \rgRotRes\,px to \rgOursRes\,px and the median number of inliers rises from \rgRotInl{} to \rgOursInl{} (\cref{tab:raw_geometry}). But the tiles are large enough for the matcher to absorb the unaligned translation, and the pose obtained from the additional inliers is only slightly more accurate. \Cref{tab:plane_vs_rot} locates the benefit: the plane-induced warp helps where a rotation leaves the largest misalignment, for oblique views (\prTiltHi{} points of \SRe{} \ci{\prTiltHiLo}{\prTiltHiHi} at tilts of $35^\circ$--$45^\circ$) and for large altitude changes (\prDzHi{} \ci{\prDzHiLo}{\prDzHiHi} above 25\,m), and brings nothing for near-nadir views or small altitude changes.

\emph{Warp under the true pose.} Sec.~4.3 states that the imagination is only as good as the hypothesis that guides it. At the accuracy of the coarse stage this is not the limiting factor: computing the warp from the ground-truth pose on \rgOracleN{} of the test pairs yields the same number of inliers (\rgOracleInl{} against \rgOursSubInl) and the same accuracy of the re-estimated pose (median direction error \rgOracleDir$^\circ$ against \rgOursSubDir$^\circ$) as the warp computed from the coarse pose (\cref{tab:raw_geometry}). The coarse pose is accurate enough for the alignment, which agrees with the finding above that a rougher alignment already suffices.

\emph{No pose-guided alignment.} To see whether the coarse pose is needed at all, we also run the tile-wise pass with no warp guided by it: the target is only rotated by the multiple of $90^\circ$ found by the rotation search of the first pass (\rgNoAlignN{} test pairs, raw geometry only, last block of \cref{tab:raw_geometry}). The pass then fails: the median number of inliers is \rgNoAlignInl, \rgNoAlignLtHundred\% of the pairs have fewer than 100, and the median direction error is $\rgNoAlignDir^\circ$, against \rgRotSubInl{} inliers and $\rgRotSubDir^\circ$ on the same pairs when the target is rotated by the heading change of the coarse pose. Without guidance, tiles at the same position of the two images do not show the same ground (median displacement \rgNoAlignRes\,px, half a tile), and the residual rotation of up to $45^\circ$ is beyond what the matcher tolerates. Matching every source tile against every target tile would remove the first obstacle at nine times the cost; we did not run it. The coarse pose is therefore necessary for the second pass, and its rotation is the part that matters most.

\emph{Warp without extra budget.} Keeping the plane-induced warp but matching a single $1280{\times}720$ tile, the budget of the zoom pass, brings \ctOneTiledSRd{} points of \SRd{} and \ctOneTiledSRe{} of \SRe{} (\ctOneTileSRd\% and \ctOneTileSRe\%), \ie \ctOneTileShareD\% and \ctOneTileShareE\% of the full gain. The median number of inliers is then \rgOneTileInl, against \rgZoomInl{} for the zoom pass and \rgOursInl{} with tiles. The remainder of the gain requires the tiles at the original resolution.

\emph{Warp at the resolution of a generative model.} Rendering the same, geometrically exact, warp at $512{\times}288$ pixels, the resolution at which a latent-diffusion model would synthesise the view, and matching it there, does not refine the coarse pose at all: \ctLowResSRd\% \SRd{} and \ctLowResSRe\% \SRe, \ie \ctLowResdSRd{} and \ctLowResdSRe{} points relative to the coarse stage. The poses re-estimated at that resolution have a median direction error of \rgLowResDir$^\circ$ from \rgLowResInl{} inliers (\cref{tab:raw_geometry}), far worse than the coarse estimate they are meant to correct. This supports with a measurement the argument of Sec.~4.3 that the imagined view must be accurate at the pixel level of the original image: no view synthesised at 512 pixels, however faithful, can refine a pose that is already accurate to a metre (\cref{sec:related}).

\noindent\textbf{(c) How the re-estimated pose is used.} Using the pose re-estimated from $\mathcal{M}_{\mathrm{im}}$ directly, in place of the coarse pose and with the coarse distance, already brings \ctDirectdSRe{} points of \SRe{} (\ctDirectShareE\% of the full gain) and leaves \SRd{} unchanged, since the length is not updated. Averaging the candidates with equal weights is harmful (\ctUniformdSRd{} points of \SRd): an unreliable candidate must be recognised as such. The predicted-error weights of Eq.~(2) alone, without the residual regressor, recover \ctLearnedNoResShareD\% and \ctLearnedNoResShareE\% of the full gain in \SRd{} and \SRe. Removing the hypothesis-verification scores changes the result by less than half a point (\ctNoVerSRd\% against \ctFullSRd\% \SRd, \ctNoVerSRe\% against \ctFullSRe\% \SRe), within the uncertainty of the comparison: they are a minor component, and Sec.~4.3 gives them more weight than they have. Keeping the scale of the coarse stage, \ie not re-triangulating the scale cues from $\mathcal{M}_{\mathrm{im}}$, removes most of the gain in \SRd{} (\ctNoScaledSRd{} instead of \ctFulldSRd{} points) and about half of the gain in \SRe{} (\ctNoScaledSRe{} instead of \ctFulldSRe).

\noindent\textbf{(d) A second match set without any imagination.} The strongest control comes from \cref{sec:siftfusion}. Giving the fusion of the coarse stage the pose and scale cues of a classical SIFT match set, computed on the half-resolution images without any pose guidance or warp, yields \sfuCoarseSiftSRd\% \SRd{} and \sfuCoarseSiftSRe\% \SRe, which is \sfuAddCoarseSRd{} and \sfuAddCoarseSRe{} points above the coarse stage, where the imagination-refined stage brings \ctFulldSRd{} and \ctFulldSRe{} on the same pairs.

\noindent\textbf{Attribution.} The controls separate the sources of the refinement gain as follows. (1) None of it is learned capacity. (2) Most of it comes from a second matching pass at the original resolution, in co-located tiles, once the coarse pose has brought the two images into rough alignment; aligning them by the coarse rotation alone is enough for about nine tenths of the gain, whereas the plane-induced warp without the tiles gives less than half of it, and at a reduced resolution none. (3) The plane-induced homography, the specific form of imagination that the main paper puts forward, adds a significant but small \ctPlaneGainSRe{} points of \SRe{} \ci{\ctPlaneGainSReLo}{\ctPlaneGainSReHi} and nothing measurable in \SRd{} (\ctPlaneGainSRd{} \ci{\ctPlaneGainSRdLo}{\ctPlaneGainSRdHi}). (4) A second match set that is accurate enough has almost the same effect when it comes from a classical detector without any warp (control (d)). (5) The new correspondences act through more accurate directions and rotations and through a re-triangulated scale, and they must be combined with the coarse estimate by predicted reliability rather than by a plain average. We therefore correct the emphasis of the main paper: the decisive ingredient of the refinement stage is pose-guided re-matching at the original resolution, and the ground-plane warp is a modest improvement of the alignment on which it relies, not the source of the gain (\cref{sec:corrections}).

\subsection{Iterating the refinement without moving}
\label{sec:iteration}

Two reviews ask what happens when the refinement is repeated on a fixed image pair, since the abstract speaks of ``iterating the refinement''. The closed loop of Sec.~5.4 does not do this (\cref{sec:internalphysical}), so we run it here on \itN{} test pairs. The estimate of one round provides the pose that guides the imagination of the next; plane, tiles, matching, pose estimation and the fitted fusion are unchanged. In protocol~A the fusion keeps the coarse estimate as a candidate, as in the first round; in protocol~B the estimate of the previous round takes its place.

\begin{table}[t]
\centering
\caption{\textbf{Repeating the refinement on a fixed image pair, without motion} (1{,}000 test pairs, camera frame). The estimate of a round guides the imagination of the next. A: the fusion keeps the coarse estimate as a candidate; B: the estimate of the previous round replaces it. $\Delta$: paired change of \SRe{} with respect to one round. Res.: median displacement between imagined and real view.}
\label{tab:iteration}
\footnotesize
\setlength{\tabcolsep}{3pt}
\resizebox{\columnwidth}{!}{%
\begin{tabular}{@{}lcccccc@{}}
\toprule
& $d$ & \SRd & EPE med. & \SRe & $\Delta$\SRe & Res. \\
& (m) & (\%) & (m) & (\%) & (points) & (px) \\
\midrule
Coarse stage & 0.93 & 73.8 & 0.82 & 60.5 & -- & -- \\
Refinement, 1 round (paper) & 0.67 & 83.1 & 0.52 & 78.4 & -- & 7.2 \\\midrule
2 rounds, A & 0.66 & 83.5 & 0.52 & 79.4 & $+1.0$ \ci{-0.3}{+2.3} & 5.0 \\
3 rounds, A & 0.67 & 83.3 & 0.51 & 79.7 & $+1.3$ \ci{-0.3}{+2.8} & 4.5 \\
2 rounds, B & 0.65 & 83.6 & 0.52 & 80.0 & $+1.6$ \ci{+0.3}{+3.0} & 5.0 \\
3 rounds, B & 0.64 & 83.9 & 0.50 & 80.3 & $+1.9$ \ci{+0.4}{+3.5} & 4.5 \\
\bottomrule
\end{tabular}}
\end{table}


\Cref{tab:iteration} shows that the guide pose does improve from round to round: the median displacement between the imagined and the real view decreases from \itOneRes{} to \itTwoARes{} and \itThreeARes\,px. The estimated pose hardly follows. After three rounds the endpoint success changes by \itThreeASRe{} points \ci{\itThreeASReLo}{\itThreeASReHi} with protocol~A and \itThreeBSRe{} points \ci{\itThreeBSReLo}{\itThreeBSReHi} with protocol~B, the second being the only change of the table whose interval excludes zero; \SRd{} changes by \itThreeASRd{} and \itThreeBSRd{} points, within its uncertainty. An internal iteration thus brings at most about two points of endpoint success. This agrees with the control that warps under the true pose (\cref{sec:controls}): once the two views are aligned well enough for the matcher, a more accurate alignment does not yield a more accurate pose. A single round is the appropriate operating point, and each further round costs as much as the refinement itself (\cref{tab:runtime}).
